\documentclass[11pt,a4paper]{article}

%================================================
% PACKAGES
%================================================
\usepackage[utf8]{inputenc}
\usepackage[T1]{fontenc}
\usepackage[french]{babel}
\usepackage[a4paper,margin=2.1cm]{geometry}

\usepackage{amsmath,amssymb,amsthm,mathtools}
\usepackage{lmodern}
\usepackage{xcolor}
\usepackage{hyperref}
\usepackage{fancyhdr}
\usepackage{lastpage}
\usepackage{enumitem}
\usepackage{tcolorbox}
\usepackage{tikz}
\usepackage{chemfig}
\usepackage{array}
\usepackage{booktabs}
\usepackage{multicol}
\usepackage{titlesec}
\usepackage{siunitx}
\usepackage{physics}
\usepackage{graphicx}

%================================================
% MISE EN PAGE
%================================================
\pagestyle{fancy}
\fancyhf{}
\fancyhead[L]{Physique-Chimie -- Spécialité}
\fancyhead[R]{Cohésion de la matière et spectroscopie IR}
\fancyfoot[C]{\thepage\ /\pageref{LastPage}}

\hypersetup{
    colorlinks=true,
    linkcolor=blue!60!black,
    urlcolor=blue!60!black,
    pdftitle={Cours - Cohésion de la matière et spectroscopie IR},
    pdfauthor={},
    pdfsubject={Physique-Chimie Spécialité}
}

\setlength{\parindent}{0pt}
\setlength{\parskip}{0.5em}

\titleformat{\section}{\Large\bfseries\color{blue!60!black}}{\thesection.}{0.6em}{}
\titleformat{\subsection}{\large\bfseries\color{blue!45!black}}{\thesubsection.}{0.6em}{}
\titleformat{\subsubsection}{\normalsize\bfseries\color{blue!30!black}}{\thesubsubsection.}{0.6em}{}

%================================================
% BOITES
%================================================
\newtcolorbox{definitionbox}{
    colback=blue!5,
    colframe=blue!60!black,
    title=Définition,
    fonttitle=\bfseries
}

\newtcolorbox{proprietebox}{
    colback=green!5,
    colframe=green!50!black,
    title=Propriété,
    fonttitle=\bfseries
}

\newtcolorbox{methodebox}{
    colback=orange!8,
    colframe=orange!75!black,
    title=Méthode,
    fonttitle=\bfseries
}

\newtcolorbox{exemplebox}{
    colback=purple!5,
    colframe=purple!60!black,
    title=Exemple,
    fonttitle=\bfseries
}

\newtcolorbox{attentionbox}{
    colback=red!5,
    colframe=red!65!black,
    title=Attention,
    fonttitle=\bfseries
}

\newtcolorbox{ideebox}{
    colback=gray!8,
    colframe=black!70,
    title=Idée physique,
    fonttitle=\bfseries
}

\newtcolorbox{analogebox}{
    colback=yellow!10,
    colframe=yellow!50!black,
    title=Analogie,
    fonttitle=\bfseries
}

\newtcolorbox{comprendrebox}{
    colback=cyan!7,
    colframe=cyan!50!black,
    title=Pour bien comprendre,
    fonttitle=\bfseries
}

\newtcolorbox{bilanbox}{
    colback=teal!6,
    colframe=teal!60!black,
    title=Bilan,
    fonttitle=\bfseries
}

%================================================
% DOCUMENT
%================================================
\begin{document}

\begin{center}
    {\LARGE \textbf{Cours de Physique-Chimie}}\\[0.3em]
    {\Large \textbf{Cohésion de la matière et spectroscopie infrarouge}}\\[0.7em]
    \rule{0.9\linewidth}{0.6pt}
\end{center}

\tableofcontents
\vspace{1em}

%================================================
\section{Introduction générale}
%================================================

Quand on regarde la matière à notre échelle, on voit des objets, des liquides, des solides, des gaz.  
Mais à l’échelle microscopique, toute matière est constituée d’\textbf{entités chimiques} :
\begin{itemize}
    \item des \textbf{atomes},
    \item des \textbf{ions},
    \item des \textbf{molécules}.
\end{itemize}

Une question essentielle apparaît alors :

\begin{center}
\textit{Pourquoi ces entités restent-elles ensemble ?}
\end{center}

Pourquoi l’eau forme-t-elle un liquide cohérent ?  
Pourquoi le sel forme-t-il un cristal solide ?  
Pourquoi l’huile ne se mélange-t-elle pas à l’eau ?  
Pourquoi certaines molécules bouillent plus vite que d’autres ?

Le premier but du chapitre est de comprendre les \textbf{forces d’interaction} responsables de la \textbf{cohésion de la matière}.  

Le second but est de comprendre comment la \textbf{spectroscopie infrarouge (IR)} permet d’identifier certaines liaisons et donc certaines fonctions chimiques dans une molécule.

\begin{ideebox}
Le grand fil directeur du chapitre est le suivant :

\begin{center}
\textbf{Structure microscopique} \(\longrightarrow\) \textbf{interactions} \(\longrightarrow\) \textbf{propriétés observables}
\end{center}

Autrement dit, ce qui se passe entre les particules explique ce qu’on voit à notre échelle.
\end{ideebox}

%================================================
\section{Cohésion de la matière}
%================================================

\subsection{Qu’appelle-t-on cohésion de la matière ?}

\begin{definitionbox}
On appelle \textbf{cohésion de la matière} l’ensemble des interactions qui maintiennent ensemble les entités microscopiques d’un corps.
\end{definitionbox}

\begin{comprendrebox}
Dire qu’une matière est ``cohésive'', cela signifie qu’il existe des interactions qui empêchent ses constituants de se disperser complètement.
\end{comprendrebox}

\begin{analogebox}
Imagine une foule :
\begin{itemize}
    \item si chaque personne est totalement indépendante, la foule se disperse ;
    \item s’il existe des liens entre les personnes, le groupe reste rassemblé.
\end{itemize}
Dans la matière, ce sont les interactions entre entités qui jouent le rôle de ces ``liens''.
\end{analogebox}

\subsection{Deux niveaux à distinguer}

Il faut distinguer :
\begin{itemize}
    \item les \textbf{liaisons fortes}, qui maintiennent ensemble les atomes dans une même entité ;
    \item les \textbf{interactions faibles}, qui agissent entre entités différentes.
\end{itemize}

\begin{attentionbox}
Il est fondamental de ne pas confondre :
\begin{itemize}
    \item ce qui tient les atomes \textbf{à l’intérieur} d’une molécule ;
    \item ce qui attire \textbf{deux molécules distinctes} l’une vers l’autre.
\end{itemize}
\end{attentionbox}

%================================================
\section{Les liaisons fortes}
%================================================

\subsection{La liaison covalente}

\begin{definitionbox}
Une \textbf{liaison covalente} est une liaison chimique dans laquelle deux atomes mettent en commun un ou plusieurs doublets d’électrons.
\end{definitionbox}

\begin{ideebox}
Chaque atome cherche une configuration électronique plus stable. Partager des électrons permet souvent d’atteindre une situation énergétique plus favorable.
\end{ideebox}

\begin{analogebox}
On peut voir la liaison covalente comme deux personnes qui tiennent une même corde : la corde représente le doublet électronique partagé.  
Aucune des deux ne possède seule toute la corde : elle est commune.
\end{analogebox}

\begin{exemplebox}
Exemples :
\begin{itemize}
    \item \(\mathrm{H_2}\) : une liaison covalente simple ;
    \item \(\mathrm{O_2}\) : une liaison covalente double ;
    \item \(\mathrm{H_2O}\) : deux liaisons covalentes O--H ;
    \item \(\mathrm{CO_2}\) : deux liaisons doubles C=O.
\end{itemize}
\end{exemplebox}

\begin{comprendrebox}
La liaison covalente est une \textbf{liaison forte}.  
Elle explique pourquoi une molécule garde sa structure.
\end{comprendrebox}

\subsection{La liaison ionique}

\begin{definitionbox}
Une \textbf{liaison ionique} résulte de l’attraction électrostatique entre deux ions de charges opposées.
\end{definitionbox}

\begin{ideebox}
Des charges opposées s’attirent : c’est une conséquence de l’interaction électrostatique.
\end{ideebox}

\begin{analogebox}
C’est un peu comme deux aimants orientés de façon à s’attirer, sauf qu’ici il s’agit de charges électriques et non de pôles magnétiques.
\end{analogebox}

\begin{exemplebox}
Le chlorure de sodium \(\mathrm{NaCl}\) est constitué d’ions :
\[
\mathrm{Na^+} \qquad \text{et} \qquad \mathrm{Cl^-}
\]
La cohésion du cristal provient de l’attraction entre ions de charges opposées.
\end{exemplebox}

\begin{comprendrebox}
Dans un solide ionique, il n’y a pas une petite ``molécule de sel'' isolée : les ions s’organisent dans tout un \textbf{réseau cristallin}.
\end{comprendrebox}

\begin{attentionbox}
Ne pas imaginer un solide ionique comme une simple juxtaposition de molécules neutres.  
Il s’agit d’un \textbf{empilement ordonné d’ions}.
\end{attentionbox}

%================================================
\section{Les interactions intermoléculaires}
%================================================

\subsection{Pourquoi des molécules neutres peuvent-elles s’attirer ?}

Deux molécules peuvent être globalement neutres et pourtant interagir.  
Cela vient de la manière dont les charges sont réparties en leur sein.

\begin{ideebox}
Une molécule neutre n’est pas forcément ``uniforme'' électriquement.  
Il peut exister des zones légèrement plus positives et des zones légèrement plus négatives.
\end{ideebox}

\subsection{Les interactions de Van der Waals}

\begin{definitionbox}
Les \textbf{interactions de Van der Waals} sont des interactions faibles entre molécules, dues à des attractions électrostatiques entre dipôles permanents, induits ou instantanés.
\end{definitionbox}

\begin{comprendrebox}
Même lorsque les molécules ne sont pas ioniques et ne portent pas de charge globale, elles peuvent se perturber mutuellement et s’attirer faiblement.
\end{comprendrebox}

\begin{analogebox}
Imagine deux nuages souples de charges qui se déforment légèrement quand ils s’approchent.  
Cette petite déformation crée une attraction faible, mais réelle.
\end{analogebox}

\begin{proprietebox}
Les interactions de Van der Waals :
\begin{itemize}
    \item existent entre toutes les molécules ;
    \item sont faibles ;
    \item deviennent plus marquées pour les grosses molécules, plus déformables.
\end{itemize}
\end{proprietebox}

\begin{exemplebox}
Le diiode \(\mathrm{I_2}\) est constitué de molécules apolaires.  
Pourtant, le solide diiode existe : sa cohésion est assurée par les interactions de Van der Waals.
\end{exemplebox}

\begin{comprendrebox}
Ces interactions expliquent notamment :
\begin{itemize}
    \item pourquoi certains gaz peuvent être liquéfiés ;
    \item pourquoi certaines molécules restent liquides à température ambiante ;
    \item pourquoi les températures d’ébullition augmentent souvent avec la taille des molécules.
\end{itemize}
\end{comprendrebox}

\subsection{La liaison hydrogène}

\begin{definitionbox}
Une \textbf{liaison hydrogène} est une interaction attractive entre :
\begin{itemize}
    \item un atome d’hydrogène déjà lié à un atome très électronégatif (\(\mathrm{O}\), \(\mathrm{N}\), \(\mathrm{F}\)) ;
    \item et un autre atome électronégatif voisin porteur d’un doublet non liant.
\end{itemize}
\end{definitionbox}

\begin{ideebox}
Quand l’hydrogène est lié à un atome très électronégatif, la liaison est fortement polarisée.  
L’hydrogène devient alors partiellement positif et peut être attiré par une zone riche en électrons d’une autre molécule.
\end{ideebox}

\begin{analogebox}
On peut voir la liaison hydrogène comme une petite ``accroche'' entre molécules, plus forte qu’une simple attraction faible ordinaire, mais bien moins forte qu’une véritable liaison covalente.
\end{analogebox}

\begin{proprietebox}
La liaison hydrogène :
\begin{itemize}
    \item est plus forte qu’une interaction de Van der Waals classique ;
    \item joue un rôle majeur dans l’eau, les alcools, certaines biomolécules ;
    \item influence fortement les températures d’ébullition et la solubilité.
\end{itemize}
\end{proprietebox}

\begin{exemplebox}
Dans l’eau, l’hydrogène d’une molécule peut interagir avec l’oxygène d’une autre molécule : il y a alors liaison hydrogène.
\end{exemplebox}

\begin{comprendrebox}
Sans les liaisons hydrogène, l’eau n’aurait pas du tout les mêmes propriétés : elle serait beaucoup moins cohésive et probablement liquide dans des conditions très différentes.
\end{comprendrebox}

%================================================
\section{Polarité des liaisons et des molécules}
%================================================

\subsection{L’électronégativité}

\begin{definitionbox}
L’\textbf{électronégativité} d’un atome mesure sa capacité à attirer vers lui les électrons engagés dans une liaison chimique.
\end{definitionbox}

\begin{ideebox}
Tous les atomes n’attirent pas les électrons avec la même intensité.
\end{ideebox}

\begin{analogebox}
Imagine une couverture tenue par deux personnes :
\begin{itemize}
    \item si elles tirent avec la même force, la couverture reste au milieu ;
    \item si l’une tire plus fort, la couverture se décale vers elle.
\end{itemize}
Dans une liaison polaire, le doublet électronique est ``tiré'' davantage par l’atome le plus électronégatif.
\end{analogebox}

\subsection{Liaison polaire}

\begin{definitionbox}
Une \textbf{liaison polaire} est une liaison covalente dans laquelle le doublet électronique est davantage attiré par l’un des deux atomes.
\end{definitionbox}

On note alors :
\[
\delta^- \quad \text{sur l’atome le plus électronégatif}, \qquad
\delta^+ \quad \text{sur l’autre}.
\]

\begin{exemplebox}
Dans \(\mathrm{H_2O}\), l’oxygène est plus électronégatif que l’hydrogène.  
Les liaisons O--H sont donc polaires.
\end{exemplebox}

\subsection{Molécule polaire ou apolaire}

\begin{definitionbox}
Une molécule est \textbf{polaire} lorsque la répartition des charges y est dissymétrique.
\end{definitionbox}

\begin{definitionbox}
Une molécule est \textbf{apolaire} lorsque la répartition des charges est globalement symétrique.
\end{definitionbox}

\begin{ideebox}
Il faut prendre en compte deux choses :
\begin{itemize}
    \item la polarité des liaisons ;
    \item la géométrie de la molécule.
\end{itemize}
\end{ideebox}

\begin{analogebox}
Même si plusieurs personnes tirent sur un objet, tout dépend de leur disposition :
\begin{itemize}
    \item si elles tirent symétriquement, les effets se compensent ;
    \item si elles tirent dans des directions non compensées, l’objet est entraîné dans une direction.
\end{itemize}
C’est exactement l’idée pour une molécule polaire.
\end{analogebox}

\begin{exemplebox}
\textbf{Exemples :}
\begin{itemize}
    \item \(\mathrm{H_2O}\) : molécule coudée, donc polaire ;
    \item \(\mathrm{CO_2}\) : liaisons polaires mais molécule linéaire et symétrique, donc apolaire ;
    \item \(\mathrm{CH_4}\) : répartition symétrique, molécule apolaire.
\end{itemize}
\end{exemplebox}

\begin{attentionbox}
Erreur très fréquente :  
\textbf{Une molécule qui possède des liaisons polaires n’est pas forcément polaire.}
\end{attentionbox}

%================================================
\section{Conséquences macroscopiques : solubilité et miscibilité}
%================================================

\subsection{Dissolution}

\begin{definitionbox}
La \textbf{dissolution} est le processus par lequel une espèce chimique se disperse dans un solvant.
\end{definitionbox}

\begin{comprendrebox}
Se dissoudre, ce n’est pas ``disparaître''.  
C’est se répartir à l’échelle microscopique parmi les molécules du solvant.
\end{comprendrebox}

\subsection{La règle ``le semblable dissout le semblable''}

\begin{proprietebox}
En première approche :
\begin{center}
\textbf{Le semblable dissout le semblable.}
\end{center}
\begin{itemize}
    \item les espèces polaires se dissolvent plutôt dans les solvants polaires ;
    \item les espèces apolaires se dissolvent plutôt dans les solvants apolaires.
\end{itemize}
\end{proprietebox}

\begin{analogebox}
Pour qu’un mélange se fasse bien, il faut que les nouvelles interactions créées soient suffisamment favorables.  
C’est un peu comme intégrer quelqu’un dans un groupe : cela marche mieux si la personne peut créer des liens du même type que ceux déjà présents.
\end{analogebox}

\begin{exemplebox}
\begin{itemize}
    \item L’eau est polaire.
    \item L’huile est globalement apolaire.
\end{itemize}
Donc :
\begin{itemize}
    \item l’eau et l’huile se mélangent mal ;
    \item l’éthanol, qui possède un groupe O--H, se mélange bien à l’eau ;
    \item beaucoup d’espèces ioniques se dissolvent dans l’eau.
\end{itemize}
\end{exemplebox}

\subsection{Dissolution d’un solide ionique dans l’eau}

L’eau étant polaire, ses molécules peuvent entourer les ions du cristal.

\begin{itemize}
    \item le côté partiellement négatif de l’eau s’oriente vers les cations ;
    \item le côté partiellement positif s’oriente vers les anions.
\end{itemize}

\begin{definitionbox}
On appelle \textbf{hydratation} l’entourage d’un ion par des molécules d’eau.
\end{definitionbox}

\begin{exemplebox}
Pour le chlorure de sodium :
\[
\mathrm{NaCl_{(s)} \rightarrow Na^+_{(aq)} + Cl^-_{(aq)}}
\]
\end{exemplebox}

\begin{comprendrebox}
La dissolution est possible si les interactions eau--ions compensent la cohésion du cristal.  
Autrement dit, l’eau doit ``arracher'' les ions au réseau cristallin.
\end{comprendrebox}

\begin{analogebox}
Imagine un mur de briques très bien assemblées.  
Pour retirer une brique, il faut qu’une autre interaction soit capable de la saisir suffisamment fort.  
Les molécules d’eau jouent ici ce rôle pour les ions.
\end{analogebox}

\subsection{Miscibilité}

\begin{definitionbox}
Deux liquides sont \textbf{miscibles} s’ils forment un mélange homogène en toute proportion.
\end{definitionbox}

\begin{exemplebox}
\begin{itemize}
    \item eau + éthanol : miscibles ;
    \item eau + huile : non miscibles.
\end{itemize}
\end{exemplebox}

\begin{attentionbox}
\textbf{Dissolution} et \textbf{miscibilité} ne désignent pas exactement la même chose :
\begin{itemize}
    \item dissolution : un soluté se disperse dans un solvant ;
    \item miscibilité : deux liquides se mélangent.
\end{itemize}
\end{attentionbox}

%================================================
\section{Le cas central de l’eau}
%================================================

\subsection{Pourquoi l’eau est-elle si particulière ?}

L’eau joue un rôle central en physique-chimie, dans les sciences du vivant, dans l’environnement et dans la vie quotidienne.

\begin{ideebox}
L’eau est une petite molécule, mais elle possède des propriétés très remarquables, liées à sa polarité et aux liaisons hydrogène.
\end{ideebox}

\subsection{Polarité de l’eau}

La molécule d’eau possède :
\begin{itemize}
    \item deux liaisons O--H polaires ;
    \item une géométrie coudée.
\end{itemize}

Les effets des deux liaisons ne se compensent donc pas : l’eau est polaire.

\subsection{Conséquences}

\begin{proprietebox}
La polarité de l’eau et les liaisons hydrogène expliquent :
\begin{itemize}
    \item sa forte cohésion ;
    \item sa tension superficielle ;
    \item sa capacité à dissoudre de nombreuses espèces ioniques ou polaires ;
    \item sa température d’ébullition relativement élevée.
\end{itemize}
\end{proprietebox}

\begin{analogebox}
Les molécules d’eau ne sont pas simplement côte à côte : elles forment un réseau d’``accroches'' temporaires entre elles grâce aux liaisons hydrogène.
\end{analogebox}

%================================================
\section{Bilan sur la cohésion de la matière}
%================================================

\begin{bilanbox}
Pour expliquer la cohésion d’une substance, il faut toujours se poser les questions suivantes :
\begin{enumerate}
    \item Quelles sont les entités présentes ?  
    ions ? molécules ? atomes ?
    \item Quels types d’interactions peuvent exister ?
    \item Ces interactions sont-elles fortes ou faibles ?
    \item Quelles conséquences a cela sur les propriétés observées ?
\end{enumerate}
\end{bilanbox}

\begin{methodebox}
Méthode de raisonnement :
\begin{enumerate}
    \item identifier la nature des entités ;
    \item repérer leur polarité éventuelle ;
    \item déterminer les interactions possibles ;
    \item relier cela à :
    \begin{itemize}
        \item l’état physique ;
        \item la solubilité ;
        \item la miscibilité ;
        \item la température d’ébullition ou de fusion.
    \end{itemize}
\end{enumerate}
\end{methodebox}

%================================================
\section{Spectroscopie infrarouge}
%================================================

\subsection{Principe général}

\begin{definitionbox}
La \textbf{spectroscopie infrarouge} est une technique d’analyse qui consiste à envoyer un rayonnement infrarouge sur une espèce chimique et à repérer quelles radiations sont absorbées.
\end{definitionbox}

\begin{ideebox}
Les atomes d’une molécule ne sont pas immobiles.  
Ils vibrent autour de leurs positions d’équilibre.
\end{ideebox}

\begin{analogebox}
On peut représenter une liaison chimique comme un petit ressort reliant deux masses.  
Ce ressort peut s’étirer, se comprimer ou se déformer.  
Le rayonnement infrarouge peut mettre en vibration ce ``ressort''.
\end{analogebox}

\begin{comprendrebox}
Chaque type de liaison a ses propres modes de vibration et absorbe donc certaines fréquences particulières du rayonnement IR.
\end{comprendrebox}

\subsection{Pourquoi cette technique est-elle utile ?}

Comme certaines liaisons absorbent dans des zones caractéristiques, le spectre IR permet d’identifier des \textbf{groupes d’atomes} et donc des \textbf{fonctions chimiques}.

\begin{proprietebox}
La spectroscopie IR permet notamment de repérer :
\begin{itemize}
    \item un groupe O--H ;
    \item une liaison C=O ;
    \item une liaison N--H ;
    \item des liaisons C--H.
\end{itemize}
\end{proprietebox}

%================================================
\section{Lecture d’un spectre IR}
%================================================

\subsection{Axes du spectre}

Un spectre IR représente généralement :
\begin{itemize}
    \item en abscisse : le \textbf{nombre d’onde} \(\sigma\), en \(\si{cm^{-1}}\) ;
    \item en ordonnée : la \textbf{transmittance} ou l’\textbf{absorbance}.
\end{itemize}

\begin{definitionbox}
Le nombre d’onde est défini par :
\[
\sigma=\frac{1}{\lambda}
\]
où \(\lambda\) est la longueur d’onde exprimée en centimètres.
\end{definitionbox}

\begin{proprietebox}
Plus le nombre d’onde est grand :
\begin{itemize}
    \item plus la longueur d’onde est petite ;
    \item plus l’énergie associée est élevée.
\end{itemize}
\end{proprietebox}

\begin{comprendrebox}
Dans les spectres usuels du lycée, les absorptions apparaissent souvent sous forme de \textbf{creux} car l’ordonnée est la transmittance.
\end{comprendrebox}

\subsection{Que signifie une bande d’absorption ?}

Une bande d’absorption signifie qu’une partie du rayonnement incident a été absorbée parce qu’elle correspond à une vibration possible d’une liaison.

\begin{analogebox}
Si l’on pousse un système périodique à la bonne fréquence, il répond fortement.  
C’est comme pousser quelqu’un sur une balançoire au bon rythme : l’oscillation devient efficace.
\end{analogebox}

%================================================
\section{Bandes caractéristiques importantes}
%================================================

\subsection{Tableau essentiel}

\begin{center}
\renewcommand{\arraystretch}{1.35}
\begin{tabular}{>{\bfseries}p{3.5cm} p{3.5cm} p{3.8cm} p{4cm}}
\toprule
Liaison / groupe & Zone typique & Aspect & Interprétation usuelle \\
\midrule
O--H (alcool) & \(3200-3600\ \si{cm^{-1}}\) & large & présence d’un alcool \\
O--H (acide) & \(2500-3200\ \si{cm^{-1}}\) & très large & acide carboxylique \\
C=O & \(1650-1750\ \si{cm^{-1}}\) & forte, nette & groupe carbonyle \\
C--H & \(2800-3100\ \si{cm^{-1}}\) & assez nette & très fréquent, peu discriminant seul \\
N--H & \(3300-3500\ \si{cm^{-1}}\) & souvent assez fine à moyenne & amine ou amide \\
\bottomrule
\end{tabular}
\end{center}

\subsection{La bande O--H}

\begin{exemplebox}
\textbf{Alcool :} bande large vers \(3200-3600\ \si{cm^{-1}}\).
\end{exemplebox}

\begin{exemplebox}
\textbf{Acide carboxylique :} bande très large vers \(2500-3200\ \si{cm^{-1}}\), souvent accompagnée d’une bande C=O.
\end{exemplebox}

\begin{comprendrebox}
La largeur de la bande O--H vient souvent du fait que les groupes O--H participent à des liaisons hydrogène, ce qui modifie les vibrations possibles.
\end{comprendrebox}

\subsection{La bande C=O}

\begin{proprietebox}
Une bande intense vers \(1700\ \si{cm^{-1}}\) est très caractéristique d’un \textbf{groupe carbonyle}.
\end{proprietebox}

\begin{exemplebox}
On trouve un groupe carbonyle dans :
\begin{itemize}
    \item les aldéhydes ;
    \item les cétones ;
    \item les acides carboxyliques ;
    \item les esters.
\end{itemize}
\end{exemplebox}

\begin{analogebox}
Une liaison double C=O est un ``ressort'' plus raide qu’une liaison simple : elle vibre donc dans une zone très caractéristique.
\end{analogebox}

%================================================
\section{Méthode d’analyse d’un spectre IR}
%================================================

\begin{methodebox}
Pour analyser un spectre IR :
\begin{enumerate}
    \item repérer les bandes les plus marquantes ;
    \item chercher d’abord un éventuel O--H ;
    \item chercher ensuite un éventuel C=O ;
    \item croiser les informations ;
    \item proposer la fonction chimique compatible.
\end{enumerate}
\end{methodebox}

\subsection{Exemple 1 : alcool}

On observe :
\begin{itemize}
    \item une large bande vers \(3300\ \si{cm^{-1}}\) ;
    \item aucune bande forte vers \(1700\ \si{cm^{-1}}\).
\end{itemize}

\textbf{Conclusion :}
\begin{itemize}
    \item présence probable d’un groupe O--H ;
    \item absence de groupe carbonyle ;
    \item on pense à un \textbf{alcool}.
\end{itemize}

\subsection{Exemple 2 : composé carbonylé}

On observe :
\begin{itemize}
    \item une bande forte vers \(1715\ \si{cm^{-1}}\) ;
    \item pas de large bande O--H.
\end{itemize}

\textbf{Conclusion :}
\begin{itemize}
    \item présence d’un groupe C=O ;
    \item pas d’alcool ni d’acide carboxylique ;
    \item on peut penser à une \textbf{cétone} ou un \textbf{aldéhyde}.
\end{itemize}

\subsection{Exemple 3 : acide carboxylique}

On observe :
\begin{itemize}
    \item une bande forte vers \(1700\ \si{cm^{-1}}\) ;
    \item une très large bande entre \(2500\) et \(3200\ \si{cm^{-1}}\).
\end{itemize}

\textbf{Conclusion :}
\begin{itemize}
    \item présence simultanée d’un groupe C=O et d’un O--H acide ;
    \item on identifie un \textbf{acide carboxylique}.
\end{itemize}

%================================================
\section{Lien entre spectroscopie IR et cohésion de la matière}
%================================================

\begin{ideebox}
Ces deux parties du cours se rejoignent :  
la structure d’une molécule détermine les interactions qu’elle peut former, et ces interactions influencent parfois le spectre IR.
\end{ideebox}

\subsection{Exemple du groupe O--H}

Si une molécule possède un groupe O--H :
\begin{itemize}
    \item elle peut souvent former des liaisons hydrogène ;
    \item cela influence sa solubilité dans l’eau ;
    \item cela influence aussi la forme de sa bande IR.
\end{itemize}

\begin{comprendrebox}
Ainsi, le spectre IR ne sert pas seulement à ``nommer une fonction'' : il donne aussi des indices sur le comportement microscopique de la molécule.
\end{comprendrebox}

%================================================
\section{Erreurs classiques à éviter}
%================================================

\subsection{Dans la cohésion de la matière}

\begin{itemize}
    \item confondre liaison covalente et interaction intermoléculaire ;
    \item croire qu’une molécule avec des liaisons polaires est forcément polaire ;
    \item oublier le rôle de la géométrie ;
    \item croire que l’eau dissout tout ;
    \item confondre dissolution et miscibilité.
\end{itemize}

\subsection{Dans la spectroscopie IR}

\begin{itemize}
    \item ne pas repérer d’abord les grandes bandes ;
    \item oublier qu’une bande vers \(1700\ \si{cm^{-1}}\) évoque fortement C=O ;
    \item conclure à partir d’une seule bande sans croiser les informations ;
    \item oublier qu’en transmittance les absorptions sont souvent des creux.
\end{itemize}

%================================================
\section{À retenir absolument}
%================================================

\begin{bilanbox}
\textbf{Cohésion de la matière}
\begin{itemize}
    \item liaison covalente : partage d’électrons ;
    \item liaison ionique : attraction entre ions de charges opposées ;
    \item interactions de Van der Waals : attractions faibles entre molécules ;
    \item liaison hydrogène : interaction forte parmi les interactions intermoléculaires usuelles ;
    \item polarité : dépend des liaisons \emph{et} de la géométrie ;
    \item solubilité : le semblable dissout le semblable.
\end{itemize}

\textbf{Spectroscopie IR}
\begin{itemize}
    \item elle met en évidence des vibrations de liaisons ;
    \item O--H alcool : \(3200-3600\ \si{cm^{-1}}\), large ;
    \item O--H acide : \(2500-3200\ \si{cm^{-1}}\), très large ;
    \item C=O : \(1650-1750\ \si{cm^{-1}}\), forte ;
    \item un spectre IR permet d’identifier certaines fonctions chimiques.
\end{itemize}
\end{bilanbox}

%================================================
\section{Mini fiche ultra-synthèse}
%================================================

\begin{center}
\renewcommand{\arraystretch}{1.4}
\begin{tabular}{>{\bfseries}p{4.1cm} p{10.3cm}}
\toprule
Notion & Sens physique essentiel \\
\midrule
Liaison covalente & Deux atomes partagent des électrons pour atteindre une situation plus stable. \\
Liaison ionique & Des ions de charges opposées s’attirent électrostatiquement. \\
Van der Waals & Même des molécules neutres peuvent s’attirer faiblement. \\
Liaison hydrogène & Interaction intermoléculaire renforcée lorsque H est lié à O, N ou F. \\
Polarité & Une molécule peut présenter un déséquilibre global de charges. \\
Solubilité & Une dissolution est favorisée si de nouvelles interactions favorables se créent. \\
Spectroscopie IR & Un rayonnement met certaines liaisons en vibration ; chaque liaison absorbe dans une zone typique. \\
Bande O--H & Large car l’environnement et les liaisons hydrogène influencent fortement la vibration. \\
Bande C=O & Très caractéristique d’un groupe carbonyle, donc très utile pour identifier une fonction. \\
\bottomrule
\end{tabular}
\end{center}

\vfill

\begin{center}
\rule{0.82\linewidth}{0.5pt}\\
\textit{Fin du cours}
\end{center}

\end{document}